\unnsection{Пример расчёта}

    В предыдущем разделе было показано, что для линейного реакционного члена
    $R(c) = -kc$ и постоянных коэффициентов точное решение уравнения
    конвекции--реакции--диффузии доступно аналитически.
    Настоящий раздел конкретизирует эту идею: формулируется задача с явно
    заданными числовыми параметрами, для неё последовательными заменами переменных
    выводится точное решение, затем численно реализуется схема Кранка--Николсон
    и проводится сравнение результатов.

    \unnsubsection{Постановка конкретной задачи}

        Рассматривается одномерное уравнение конвекции--реакции--диффузии:
        \[
            \frac{\partial c}{\partial t}
            + v\,\frac{\partial c}{\partial x}
            = D\,\frac{\partial^2 c}{\partial x^2} - k c,
            \qquad x \in \mathbb{R},\quad t > 0,
        \]
        с начальным условием в виде гауссова пика, сосредоточённого в точке $x = 0$:
        \[
            c(x, 0) = \exp\!\left(-x^2\right).
        \]
        Физически задача описывает следующую ситуацию: в начальный момент времени
        в среде, движущейся с постоянной скоростью $v$, образовалось облако примеси
        с гауссовым профилем концентрации.
        При $t > 0$ облако одновременно переносится потоком (конвекция),
        расплывается за счёт молекулярного перемешивания (диффузия)
        и постепенно убывает вследствие необратимого распада примеси по реакции первого
        порядка (например, радиоактивный распад или нейтрализация загрязнителя).
        Значения параметров задачи:
        \[
            v = 0{,}5 \ \text{м/с}, \quad
            D = 0{,}1 \ \text{м}^2\!/\text{с}, \quad
            k = 0{,}2 \ \text{с}^{-1}.
        \]
        При характерном пространственном масштабе задачи $L \sim 1$\,м глобальное
        число Пекле $\mathrm{Pe} = vL/D = 5$: конвекция заметно доминирует над
        диффузией, однако оба механизма вносят ощутимый вклад.

    \unnsubsection{Аналитическое решение}

        Точная формула получается двумя последовательными заменами переменных,
        каждая из которых устраняет один из двух «возмущающих» членов.

        \textbf{Шаг~1. Устранение реакционного слагаемого.}
        Полагаем $c(x, t) = e^{-kt}\,w(x, t)$.
        Вычислим производные и подставим в уравнение:
        \[
            e^{-kt}\!\left(\frac{\partial w}{\partial t}
            + v\,\frac{\partial w}{\partial x}\right)
            = D\,e^{-kt}\frac{\partial^2 w}{\partial x^2}.
        \]
        Сокращая на $e^{-kt} > 0$, получаем уравнение конвекции--диффузии
        без реакции:
        \[
            \frac{\partial w}{\partial t}
            + v\,\frac{\partial w}{\partial x}
            = D\,\frac{\partial^2 w}{\partial x^2}, \qquad
            w(x, 0) = e^{-x^2}.
        \]

        \textbf{Шаг~2. Устранение конвективного слагаемого.}
        Переходим в систему координат, движущуюся со скоростью $v$:
        $\xi = x - vt$, $\tau = t$.
        В новых переменных полная производная по времени при постоянном $\xi$
        равна $\partial w/\partial \tau = \partial w / \partial t + v\,\partial w / \partial x$,
        поэтому уравнение конвекции--диффузии принимает вид чистого уравнения теплопроводности:
        \[
            \frac{\partial w}{\partial \tau}
            = D\,\frac{\partial^2 w}{\partial \xi^2}, \qquad
            w(\xi, 0) = e^{-\xi^2}.
        \]
        Его решение задаётся свёрткой начального условия с фундаментальным решением:
        \[
            w(\xi, \tau)
            = \int_{-\infty}^{+\infty}
              \frac{1}{\sqrt{4\pi D\tau}}\,
              e^{-(\xi-\eta)^2/(4D\tau)}\cdot
              e^{-\eta^2}\,d\eta.
        \]
        Вычислим этот интеграл, выделив полный квадрат в показателе степени:
        \begin{align*}
            -\frac{(\xi-\eta)^2}{4D\tau} - \eta^2
            &= -\frac{1+4D\tau}{4D\tau}
              \left(\eta - \frac{\xi}{1+4D\tau}\right)^{\!2}
              - \frac{\xi^2}{1+4D\tau}.
        \end{align*}
        Интеграл от гауссовой экспоненты по $\eta$ вычисляется стандартно:
        $\int_{-\infty}^{+\infty} e^{-A(\eta - B)^2}\,d\eta = \sqrt{\pi/A}$,
        откуда
        \[
            w(\xi, \tau)
            = \frac{1}{\sqrt{1+4D\tau}}\,
              \exp\!\left(-\frac{\xi^2}{1+4D\tau}\right).
        \]
        Возвращаясь к исходным переменным ($\xi = x - vt$, $\tau = t$) и
        подставляя множитель $e^{-kt}$, приходим к \textit{точному решению
        исходной задачи}:
        \[
            \boxed{
                c(x, t)
                = \frac{e^{-kt}}{\sqrt{1+4Dt}}\,
                  \exp\!\!\left(-\frac{(x - vt)^2}{1+4Dt}\right).
            }
        \]
        Формула явно разделяет вклады трёх физических механизмов.
        \textit{Конвекция} смещает центр пика со скоростью $v$: в показателе
        степени стоит $(x - vt)$, а не $x$.
        \textit{Диффузия} расширяет пик (знаменатель $1 + 4Dt$ монотонно растёт)
        и снижает его высоту (предэкспоненциальный множитель $1/\sqrt{1+4Dt}$
        убывает обратно пропорционально ширине).
        \textit{Реакция} равномерно подавляет всё поле концентрации на $e^{-kt}$,
        не меняя форму профиля.

        В таблице~\ref{tab:cmax} приведены аналитические значения максимальной
        концентрации (достигается в центре пика $x = vt$) в контрольные моменты
        времени, а также раздельный вклад диффузии и реакции.

        \begin{table}[H]
            \centering
            \label{tab:cmax}
            \begin{tabular}{lrrrrr}
                \toprule
                & $t = 0$ & $t = 1$ & $t = 2$ & $t = 3$ & $t = 5$ \\
                \midrule
                Только диффузия ($k=0$): $1/\sqrt{1+4Dt}$
                    & 1{,}000 & 0{,}845 & 0{,}745 & 0{,}674 & 0{,}577 \\
                Только реакция ($D=0$): $e^{-kt}$
                    & 1{,}000 & 0{,}819 & 0{,}670 & 0{,}549 & 0{,}368 \\
                Оба механизма: $e^{-kt}/\sqrt{1+4Dt}$
                    & 1{,}000 & 0{,}692 & 0{,}499 & 0{,}370 & 0{,}212 \\
                \bottomrule
            \end{tabular}
            \caption{Максимальная концентрация и вклад отдельных механизмов}
        \end{table}

        Из таблицы видно, что совместное действие диффузии и реакции приводит
        к значительно более быстрому убыванию максимума, чем каждый механизм
        по отдельности: к моменту $t = 5$\,с максимальная концентрация составляет
        лишь $21{,}2\,\%$ от начальной.

    \unnsubsection{Численный метод}

        Для верификации аналитического решения задача решается численно.
        Бесконечная прямая заменяется конечным отрезком $[x_l, x_r] = [-2, 10]$\,м
        с однородными граничными условиями Дирихле $c(x_l, t) = c(x_r, t) = 0$.
        Границы выбраны достаточно удалёнными: на всём промежутке $t \in [0, 5]$\,с
        точное решение в точках $x = -2$ и $x = 10$ не превышает $3\cdot10^{-4}$,
        поэтому замена бесконечной области конечной практически не влияет на решение.

        Введём равномерные сетки: $x_i = x_l + i\,\Delta x$
        ($i = 0, \ldots, N$,\ $\Delta x = 12/N$)
        и $t^n = n\,\Delta t$ ($n = 0, \ldots, M$,\ $\Delta t = T/M$).
        Схема Кранка--Николсон усредняет пространственный оператор между
        временными слоями $n$ и $n+1$:
        \[
            \frac{c_i^{n+1} - c_i^n}{\Delta t}
            = \frac{1}{2}\left[\mathcal{L}(c^n)_i
                              + \mathcal{L}(c^{n+1})_i\right],
        \]
        где пространственный оператор аппроксимируется центральными разностями:
        \[
            \mathcal{L}(c)_i
            = D\,\frac{c_{i+1} - 2c_i + c_{i-1}}{\Delta x^2}
            - v\,\frac{c_{i+1} - c_{i-1}}{2\,\Delta x}
            - k\,c_i.
        \]
        Применение центральных разностей для конвективного члена оправдано
        малостью ячеечного числа Пекле:
        \[
            \mathrm{Pe}_h
            = \frac{v\,\Delta x}{D}
            = \frac{0{,}5 \cdot 0{,}02}{0{,}1}
            = 0{,}1 \ll 2.
        \]
        Введём безразмерные коэффициенты:
        \[
            \alpha = \frac{D\,\Delta t}{2\,\Delta x^2}, \qquad
            \beta  = \frac{v\,\Delta t}{4\,\Delta x},   \qquad
            \gamma = \frac{k\,\Delta t}{2}.
        \]
        Перенося неявную часть в левую сторону, на каждом временном шаге
        необходимо решить \textit{трёхдиагональную} систему
        $A\,\mathbf{c}^{n+1} = B\,\mathbf{c}^n$
        для вектора значений во внутренних узлах.
        Матрицы $A$ (неявная) и $B$ (явная) имеют структуру:
        \[
            A:\;
            \begin{cases}
                \text{поддиагональ:} & -(\alpha+\beta), \\
                \text{главная:}      & \phantom{-}1 + 2\alpha + \gamma, \\
                \text{наддиагональ:} & -(\alpha-\beta);
            \end{cases}
            \qquad
            B:\;
            \begin{cases}
                \text{поддиагональ:} & \alpha+\beta, \\
                \text{главная:}      & 1 - 2\alpha - \gamma, \\
                \text{наддиагональ:} & \alpha-\beta.
            \end{cases}
        \]
        Схема Кранка--Николсон является \textit{безусловно устойчивой} (для
        диффузии--реакции) и имеет \textit{второй порядок точности} по времени
        и по пространству: $O(\Delta t^2 + \Delta x^2)$.
        Трёхдиагональная система решается алгоритмом прогонки,
        реализованным в функции \texttt{scipy.linalg.solve\_banded}.

    \unnsubsection{Реализация на Python}

        Ниже приведена полная программа.
        Она вычисляет аналитическое решение по явной формуле,
        параллельно проводит численную интеграцию схемой Кранка--Николсон,
        строит два графика и выводит таблицу погрешностей.
        Программу следует запустить до компиляции \LaTeX{}, чтобы
        сгенерировать файл \texttt{figures/crd\_solution.png}.

\begin{minted}{python}
import numpy as np
import matplotlib.pyplot as plt
from scipy.linalg import solve_banded

# Параметры задачи
v = 0.5 # скорость конвекции, м/с
D = 0.1 # коэффициент диффузии, м²/с
k = 0.2 # константа скорости реакции первого порядка, 1/с

x_left  = -2.0 # левая граница области, м
x_right = 10.0 # правая граница области, м
T       = 5.0  # конечный момент времени, с

# Параметры сетки
N  = 600                    # число пространственных интервалов
M  = 2000                   # число временных шагов
dx = (x_right - x_left) / N # шаг по пространству: 0.02 м
dt = T / M                  # шаг по времени: 0.0025 с

# Массивы узлов
x = np.linspace(x_left, x_right, N + 1)

# Начальное условие
c = np.exp(-x**2)

# Аналитическое решение
def c_exact(x, t):
    if t == 0.0:
        return np.exp(-x**2)
    denom = 1.0 + 4.0 * D * t
    return np.exp(-k * t) / np.sqrt(denom) * np.exp(-(x - v * t)**2 / denom)

# Коэффициенты схемы Кранка–Николсон
alpha = D * dt / (2.0 * dx**2) # диффузионный параметр
beta  = v * dt / (4.0 * dx)    # конвективный параметр
gamma = k * dt / 2.0           # реакционный параметр

# Матрица A
#   ab[0, j] = элемент наддиагонали  A[j-1, j] = -(alpha - beta)
#   ab[1, j] = диагональный элемент  A[j, j]   = 1 + 2*alpha + gamma
#   ab[2, j] = элемент поддиагонали  A[j+1, j] = -(alpha + beta)
n_inner = N - 1                   # число внутренних узлов
ab = np.zeros((3, n_inner))
ab[0, 1:]  = -(alpha - beta)      # наддиагональ
ab[1, :]   =  1 + 2*alpha + gamma # главная диагональ
ab[2, :-1] = -(alpha + beta)      # поддиагональ

# Временная интеграция
plot_times = [0.0, 1.0, 2.0, 3.0, 5.0] # моменты для записи профилей
snapshots  = {}   # словарь: t -> (c_численное, c_аналитическое)

def record(t_val, c_num):
    snapshots[t_val] = (c_num.copy(), c_exact(x, t_val))

record(0.0, c)

for n in range(M):
    ci = c[1:-1]   # значения во внутренних узлах на слое n

    # Правая часть системы: вектор B * c^n
    # B: поддиагональ = (alpha+beta),
    #    главная      = (1 - 2*alpha - gamma),
    #    наддиагональ = (alpha - beta)
    rhs = np.zeros(n_inner)

    # Первый внутренний узел:
    rhs[0] = ((1 - 2*alpha - gamma) * ci[0]
              + (alpha - beta) * ci[1])

    # Внутренние узлы:
    rhs[1:-1] = ((alpha + beta) * ci[:-2]
                 + (1 - 2*alpha - gamma) * ci[1:-1]
                 + (alpha - beta) * ci[2:])

    # Последний внутренний узел:
    rhs[-1] = ((alpha + beta) * ci[-2]
               + (1 - 2*alpha - gamma) * ci[-1])

    # Решение системы
    c_new_inner = solve_banded((1, 1), ab, rhs)

    # Сборка полного вектора с граничными условиями c = 0
    c_new = np.zeros(N + 1)
    c_new[1:-1] = c_new_inner
    c = c_new

    # Запись среза при достижении контрольного момента
    t_cur = (n + 1) * dt
    for pt in plot_times:
        if pt not in snapshots and abs(t_cur - pt) < dt / 2:
            record(pt, c)

# Визуализация результатов
fig, axes = plt.subplots(1, 2, figsize=(14, 5))
fig.suptitle("Решение уравнения конвекции–реакции–диффузии", fontsize=13)

# Левый график: профили концентрации в разные моменты времени
ax = axes[0]
colors = plt.cm.plasma(np.linspace(0.1, 0.85, len(plot_times)))
for i, pt in enumerate(plot_times):
    c_num, c_an = snapshots[pt]
    ax.plot(x, c_an,  color=colors[i], lw=2.0,
            label=f"t = {pt} с (аналит.)")
    ax.plot(x, c_num, color=colors[i], lw=1.5, ls="--",
            label=f"t = {pt} с (числен.)")
ax.set_xlabel("x, м")
ax.set_ylabel("c(x, t)")
ax.set_title("Профили концентрации примеси")
ax.legend(fontsize=7.5, ncol=2)
ax.set_xlim(x_left, x_right)
ax.set_ylim(-0.05, 1.1)
ax.grid(True, alpha=0.3)

# Правый график: убывание максимальной концентрации
ax2 = axes[1]
t_arr = np.linspace(0.0, T, 500)
c_max_both = np.exp(-k * t_arr) / np.sqrt(1.0 + 4.0 * D * t_arr)
c_max_diff = 1.0 / np.sqrt(1.0 + 4.0 * D * t_arr) # только диффузия (k=0)
c_max_rxn  = np.exp(-k * t_arr)                   # только реакция (D=0)
ax2.plot(t_arr, c_max_both, "k-",  lw=2.0,
         label="Оба эффекта (диффузия + реакция)")
ax2.plot(t_arr, c_max_diff, "b--", lw=1.5,
         label="Только диффузия ($k=0$)")
ax2.plot(t_arr, c_max_rxn,  "r:",  lw=1.5,
         label="Только реакция ($D=0$)")
ax2.set_xlabel("t, с")
ax2.set_ylabel("max c(x, t)")
ax2.set_title("Убывание максимальной концентрации")
ax2.legend()
ax2.grid(True, alpha=0.3)

plt.tight_layout()
plt.savefig("figures/crd_solution.png", dpi=150, bbox_inches="tight")
plt.show()

# Таблица максимальных абсолютных погрешностей
print(f"\n{"t, с":>6} | {"max |c_числ - c_аналит|":>25}")
print("-" * 35)
for pt in plot_times[1:]:
    c_num, c_an = snapshots[pt]
    err = np.max(np.abs(c_num - c_an))
    print(f"{pt:>6.1f} | {err:>25.2e}")
\end{minted}

    \unnsubsection{Результаты и их анализ}

        После запуска программы в консоли выводится таблица максимальных
        абсолютных погрешностей, а в файл сохраняются два графика
        (рис.~\ref{fig:crd_solution}).

\begin{minted}[linenos=false]{text}
  t, с |   max |c_числ - c_аналит|
-----------------------------------
   1.0 |                  8.66e-03
   2.0 |                  5.40e-03
   3.0 |                  3.72e-03
   5.0 |                  1.98e-03
\end{minted}

        \begin{figure}[H]
            \centering
            \includegraphics[width=\linewidth]{crd_solution}
            \caption{Слева: профили концентрации $c(x,t)$ в моменты $t = 0, 1, 2, 3, 5$\,с;
                     сплошные линии~--- аналитическое решение,
                     штриховые~--- схема Кранка--Николсон.
                     Справа: убывание максимальной концентрации с разбивкой
                     на вклад диффузии и реакции.}
            \label{fig:crd_solution}
        \end{figure}

        \textbf{Левый график} показывает эволюцию профиля концентрации.
        Центр гауссова пика движется вправо со скоростью $v = 0{,}5$\,м/с
        (за 5\,с смещается на 2{,}5\,м), пик постепенно уширяется и
        одновременно убывает по амплитуде.
        Штриховые (численные) кривые практически сливаются со сплошными
        (аналитическими), что свидетельствует о высокой точности схемы.

        \textbf{Правый график} наглядно демонстрирует, что суммарное убывание
        максимума (чёрная кривая) происходит значительно быстрее, чем
        от каждого механизма в отдельности:
        к моменту $t = 5$\,с только диффузия снизила бы максимум до
        $57{,}7\,\%$ от начального значения, только реакция~--- до $36{,}8\,\%$,
        а оба эффекта вместе~--- до $21{,}2\,\%$.
        Это объясняется тем, что в аналитической формуле диффузионный множитель
        $1/\sqrt{1+4Dt}$ и реакционный $e^{-kt}$ входят \textit{произведением},
        а не суммой: их совместное действие мультипликативно.

        \textbf{Таблица погрешностей}, выводимая программой в консоль,
        показывает, что максимальное отклонение численного решения от
        аналитического не превышает $10^{-2}$ и монотонно убывает со временем.
        Доминирующий источник погрешности~--- не усечение разностной схемы,
        а несоответствие граничных условий: аналитическое решение получено для
        бесконечной прямой, тогда как численно задача решается на отрезке $[-2, 10]$\,м
        с условием $c = 0$ на краях.
        В начальный момент это условие нарушено на величину
        $c(-2, 0) = e^{-4} \approx 0{,}018$; именно этим объясняется
        начальная погрешность порядка $10^{-2}$.
        По мере того как плюм смещается вправо, концентрация у левой границы
        экспоненциально убывает, что и приводит к уменьшению погрешности
        с ростом $t$.
